\section*{Problem 1}

已知矩阵 \[
A = \begin{bmatrix}
1 & 2 & 3 \\
2 & 4 & 6 \\
1 & 2 & 3 \\
\end{bmatrix}.
\]
该矩阵的各行之间存在明显的倍数关系，因此可以用低秩形式表示。

\begin{enumerate}
    \item 将矩阵 $A$ 写成外积形式 \[A = \boldsymbol{u}\boldsymbol{v}^\top.\]
    \item 求矩阵 $A$ 的秩。
    \item 若直接存储矩阵 $A$ 需要存储全部 $3 \times 3$ 个元素，而采用外积形式时只需存储向量 $\boldsymbol{u}$ 和 $\boldsymbol{v}$ 的全部分量，求：
    \begin{enumerate}
        \item 原矩阵的存储量；
        \item 外积表示下的存储量；
        \item 传输效率提升倍数（定义为压缩前存储量 / 压缩后存储量）。
    \end{enumerate}
\end{enumerate}

\section*{Problem 2}

设对称矩阵 \[
A = \begin{bmatrix}
    3 & 1 \\
    1 & 3 \\
\end{bmatrix}.
\]
利用特征值分解研究该矩阵的低秩近似。

\begin{enumerate}
    \item 求矩阵 $A$ 的全部特征值及对应的单位特征向量；
    \item 写出矩阵 $A$ 的特征值分解形式；
    \item 取最大特征值所对应的一项，写出矩阵 $A$ 的最佳秩 $1$ 近似矩阵 $\hat{A}$。
\end{enumerate}

\section*{Problem 3}

设对称矩阵 \[
A = \begin{bmatrix}
    2 & 1 & 0 \\
    1 & 2 & 0 \\
    0 & 0 & 1 \\
\end{bmatrix}
\]
\begin{enumerate}
    \item 求矩阵 $A$ 的全部特征值及对应的一组单位特征向量；
    \item 按特征值从大到小排序，并写出前三个主成分方向；
    \item 设 $\hat{A}_k$ 为保留前 $k$ 个主成分得到的近似矩阵，判断当误差要求 \[\left\|A - \hat{A}_k\right\| \le 1\] 时，至少需要保留几个主成分。
\end{enumerate}
